\section{粒子滤波用带权样本追踪非线性状态}
\label{sec:13-Machine-Learning/08-Sampling-and-Inference/05}

港口的自动导引车走到一个丁字路口，激光雷达给出的距离读数同时与向左和向右两条路径吻合。你的定位模块用一个均值加一个协方差表示当前位置，于是它给出的答案落在两条路中间那堵墙里——那个位置的真实概率接近零，却是两个峰的加权平均。这不是滤波器没调好。\hyperref[sec:13-Machine-Learning/07-Graphical-and-Sequential-Models/03]{``Kalman filter 是动态系统上的 Gaussian 条件化''} 之所以能用一对矩精确表示滤波后验，靠的是线性变换与 Gaussian 条件化让整个族保持闭合；一旦后验分成两个峰，这个表示从原理上就装不下。

粒子滤波不再预设后验单峰，而用一组带权样本表示它：

\[
p(x_t\given y_{1:t})
\approx\sum_{i=1}^{N}w_t^{(i)}\,\delta_{x_t^{(i)}}(x_t),
\qquad\sum_{i=1}^{N}w_t^{(i)}=1.
\]

每个粒子是一个当前状态的假设，权重是它与已有观测的相容程度。丁字路口的两条路各分到一批粒子，谁也不必去中间那堵墙里。换来的表示能力有价：整套方法是上一节之外的另一种 Monte Carlo，误差按样本量走，而权重退化会随时间累积。

\subsection{预测与校正的结构与 Kalman 完全同源}

一般状态空间模型只需要两个条件分布：

\[
X_t\sim p(x_t\given x_{t-1}),\qquad Y_t\sim p(y_t\given x_t).
\]

由 Markov 性，先做预测，

\[
p(x_t\given y_{1:t-1})=\int p(x_t\given x_{t-1})\,p(x_{t-1}\given y_{1:t-1})\,dx_{t-1},
\]

新观测到达后再校正，

\[
p(x_t\given y_{1:t})\propto p(y_t\given x_t)\,p(x_t\given y_{1:t-1}).
\]

这两式与 Kalman filter 一字不差，隐 Markov 模型的前向递推（见 \hyperref[sec:06-Random-Processes/04-Applications/04]{``Hidden Markov Model：从观测反推状态''}）也是同一个结构。分歧不在 Bayesian 公式，而在分布用什么表示、积分怎么近似：Kalman 用一对矩加解析积分，前向算法在有限状态上做精确求和，粒子滤波用上一时刻的加权经验分布代替预测积分，再用似然重新分配质量。校正式右端的归一化常数依旧算不出，依旧被权重的归一化悄悄消掉——这是本单元第三次看到同一招。

\subsection{权重沿时间相乘，所以它注定会塌缩}

沿着提议分布

\[
q(x_{0:t}\given y_{1:t})=q(x_{0:t-1}\given y_{1:t-1})\,q(x_t\given x_{0:t-1},y_{1:t})
\]

逐步延长每条轨迹，未归一化权重满足递推

\[
\widetilde w_t^{(i)}
=w_{t-1}^{(i)}\,
\frac{p(y_t\given x_t^{(i)})\,p(x_t^{(i)}\given x_{t-1}^{(i)})}
{q(x_t^{(i)}\given x_{0:t-1}^{(i)},y_{1:t})},
\qquad
w_t^{(i)}=\frac{\widetilde w_t^{(i)}}{\sum_j\widetilde w_t^{(j)}}.
\]

最省事的 bootstrap 粒子滤波取转移先验作提议，\(q(x_t\given x_{t-1},y_t)=p(x_t\given x_{t-1})\)，转移项上下抵消，一步只剩两件事：按动力学把粒子推一步，再按 \(p(y_t\given x_t^{(i)})\) 加权。实现是简单了，可这个提议完全没用上当前观测。雷达精度高而里程计噪声大时，推出去的粒子大多落在似然近乎为零的地方，算力在看见观测之前就已经花掉。

\begin{center}
\begin{tikzpicture}[>=stealth]
\draw[->] (-3.6,0.2) -- (3.9,0.2) node[right]{\(x_t\)};
\draw[thick] plot[smooth] coordinates
  {(0.2,0.24) (0.7,0.42) (1.0,0.90) (1.25,1.50) (1.4,1.70)
   (1.55,1.50) (1.8,0.90) (2.1,0.42) (2.6,0.24)};
\node at (3.0,1.35) {\(p(y_t\given x_t)\)};
\node[left] at (-3.5,-0.55) {预测后};
\foreach \p in {-3.0,-2.4,-1.8,-1.1,-0.5,0.1,0.7,1.3,1.9,2.6}
  \fill (\p,-0.55) circle (2pt);
\node[left] at (-3.5,-1.5) {加权后};
\foreach \p/\r in {-3.0/0.4,-2.4/0.4,-1.8/0.4,-1.1/0.4,-0.5/0.5,
                   0.1/0.7,0.7/1.6,1.3/5.2,1.9/2.1,2.6/0.6}
  \fill (\p,-1.5) circle (\r pt);
\node[left] at (-3.5,-2.45) {重采样后};
\foreach \p in {0.72,1.16,1.24,1.28,1.30,1.33,1.36,1.42,1.88,1.94}
  \fill (\p,-2.45) circle (2pt);
\draw[gray,dashed] (1.4,-2.75) -- (1.4,1.85);
\node[gray,right] at (1.95,-2.9) {\(N_{\mathrm{eff}}\) 从 \(10\) 掉到不足 \(2\)};
\end{tikzpicture}
\end{center}

\emph{圆点面积正比于权重：一次苛刻的观测就能把全部质量压到一两个粒子上。}

图中第二行画的就是权重塌缩的一步。哪怕每一步的不均衡都不算严重，权重是连乘的，几十步之后少数轨迹会吞掉几乎全部质量。可观察的诊断量与上一节的重要性采样一模一样，

\[
N_{\mathrm{eff}}=\frac{1}{\sum_{i=1}^{N}\bigl(w_t^{(i)}\bigr)^2}\in[1,N],
\]

它是权重集中度的启发式度量，不是独立样本数的严格估计，但作为报警器足够灵敏：\(N_{\mathrm{eff}}\) 掉到 \(N/10\) 以下，这一步的估计实际上由个位数的粒子撑着。

重采样是止损手段。按当前权重抽出 \(N\) 个祖先，权重重置为 \(1/N\)，高权重粒子被复制，低权重粒子消失。多项式重采样最直接，系统重采样与分层重采样引入的额外方差更小。要看清它解决了什么、又制造了什么：它治的是``质量集中在少数粒子上''，换来的新病是``大量粒子共享同一个祖先''。图中第三行十个点看着回到了等权，其实只来自三个不同的祖先。当前时刻的状态估计还站得住，整条路径的后验已经严重贫化——做长期平滑时，早期时刻往往只剩一两条祖先轨迹存活。所以通常按 \(N_{\mathrm{eff}}\) 阈值自适应地重采样而不是每步都做，需要轨迹时另配后向模拟、祖先采样或专门的粒子平滑。

一轮 bootstrap 滤波因此定型为：先按转移分布推进每个粒子，再算对数权重 \(\log\widetilde w_t^{(i)}=\log p(y_t\given x_t^{(i)})\)，用 log-sum-exp 归一化后求出所需期望 \(\sum_iw_t^{(i)}h(x_t^{(i)})\)，最后检查 \(N_{\mathrm{eff}}\) 并决定是否重采样。在对数域上做归一化是数值上的硬要求：长序列里似然连乘极易下溢成零。不过减去最大值再指数化只能防下溢，它救不了提议与后验严重错位——所有对数权重都极小时，数值是干净的，信息量是零。

顺带一提，逐步的归一化常数乘起来给出边缘似然的无偏估计，粒子方法因此能嵌进参数比较或 particle MCMC。要留心的是无偏不等于低方差：同一批有限粒子给出的似然估计可能波动极大，放进外层算法后混合会慢得离谱。这与第一节强调的那点是同一件事。

\subsection{加粒子修不了模型错设}

粒子滤波适合非线性、非 Gaussian 乃至多峰的低维到中等维状态。维数一高，观测似然会把后验质量压进一个极小的区域，落在里面的粒子比例随维数指数下降，所需粒子数随之爆炸。这是重要性采样的维数灾难在时间上的逐步累积，靠并行更多粒子并不总能扛过去。

诊断时至少要把三类失败分开。表示失败是粒子太少，稀有但要紧的模态从来没被覆盖，加粒子有效。提议失败是粒子来自转移先验、与当前观测几乎不相容，此时该改提议——用观测引导的提议、辅助粒子滤波，加粒子只是线性地缓解一个指数级的问题。模型失败是真实的转移或观测噪声压根不在模型里，粒子再多也只是更精确地逼近一个错误的后验。

区分它们的手法是变量控制：在同一段数据上改变粒子数与重采样阈值，看估计是否稳定；监控 \(N_{\mathrm{eff}}\)、独特祖先数以及多次独立运行之间的散布；把模型退化成线性 Gaussian 的版本，与 Kalman filter 的结果对照。粒子数增加而结果不收敛，先查提议与实现；结果稳定收敛却持续偏离真实轨迹，那就得回到状态空间模型本身。这几节里那个未知的归一化常数一直被巧妙地约掉，下一节要正面处理它——有些模型不给你这个机会。
